针对示向误差有界、有效接收半径未知且干扰源数量未知条件下的快速定位清除问题，本文建立“可行域交会--最小包围圆清除--有限覆盖搜索”的统一模型。对问题一，将每次示向度表示为两条半平面约束，利用线性规划判定交集为空、有界或无界，并由有界凸多边形顶点求区域直径。构造边长为 $40\unit{m}$ 的等边三角形定位区域，其最小包围圆半径为 $23.094\unit{m}>20\unit{m}$，说明以区域直径为直径的圆不能保证覆盖定位区域。

对问题二，首次获得示向度后，信号源位于半径 $1500\unit{m}$、半角 $1^\circ$ 的扇形内。为保证第二点在最小有效接收半径 $1000\unit{m}$ 下仍能收到信号，将候选域定义为首测点和扇形远端两端点处三个 $1000\unit{m}$ 圆盘的交集。给出闭式对称推荐点：在局部坐标中为 $(761.429,\pm648.248)\unit{m}$。其到全部可能源位置的最大距离恰为 $1000\unit{m}$，在 $321\times401$ 边界网格上最小交会角为 $39.592^\circ$。

对问题三、四，以边长 $600\unit{m}$ 的方格构造有限覆盖点集，并在每点扫描尚未清除的频道。全向源由 $848.53\unit{m}<1000\unit{m}$ 的距离界保证被发现；对定向源，利用源点位于方格四顶点凸包内，证明任一过源半平面至少包含一个顶点。发现信号后持续相交示向楔形，最小包围圆半径不超过 $19.5\unit{m}$ 时清除，否则进入有限的 $20\unit{m}$ 网格兜底。六轮官方演练共清除 $82/82$ 个源，请求拒绝和清除失败均为零；同种子 100 局合成回归全部清除。正式测试尚未启动，论文中的正式结果表将在真实测试后据原始日志填写。

\keywords{无线电测向；交会定位；最小包围圆；覆盖路径；稳健定位}
